\section{感悟}
通过第三周的实验练习，形成以下认识和体会：
\begin{itemize}
\item 打包与分发的重要性：借助 pyproject.toml 与 python -m build，可以把一个最小命令行包标准化地构建为 wheel，并在全新虚拟环境中仅凭该 wheel 完成安装与运行，保证了交付的可复现性。
\item 测试驱动的智能体协作：先写失败测试、再让智能体以测试通过为目标修改实现，能够形成可验证的闭环；但智能体的改动仍需人工用 git diff 审查，防止引入无关修改。
\item 协作材料的可执行性：一份好的 Issue、提交信息与评审意见应当给出环境、复现步骤、期望与实际结果，并标注风险级别，才能让协作者无需反复追问即可行动。
\item PyTorch 训练循环的闭环：一次完整的参数更新包含 zero\_grad、backward、step 三步；梯度若不清空会逐次累加，理解这一点是正确使用自动求导的前提。
\end{itemize}

从把代码打包分发，到与智能体协作，再到撰写高质量协作材料和使用深度学习框架，第三周的内容从“代码如何交付”延伸到了“如何与他人、与工具高效协作”，体现了工程实践不止于写出能运行的代码，更在于让代码与协作过程都可复现、可验证、可维护。


\subsection*{Github提交}
此处上传Github提交记录。(报告于完成后上传)
\begin{figure}[H]
    \centering
    \includegraphics[width=0.95\linewidth]{git}
\end{figure}
